\section{Original AC--MOT}
\subsection{Pipeline and information flow}
The original controller is placed between the image stream and an otherwise fixed detector--tracker pair. At an analysis step it receives inexpensive image statistics and information about the boxes retained by the preceding tracking step. It computes a scene state, maps that state to detector confidence, NMS, and resolution, executes the detector, and passes the resulting boxes to the tracker. The tracker output is stored for a later analysis step. This ordering is important: the controller does not use ground truth, future frames, or the tracker output produced by the detector action it is currently choosing.

\subsection{Five scene cues}
Let $N_t$ denote the number of boxes available from the preceding frame and let $a_{i,t}$ denote their source-image areas. The crowd cue is
\begin{equation}
C_t=\min\left(\frac{N_t}{30},1\right).
\end{equation}
The tiny-object cue is the fraction of available boxes below the historical $32\times32$ pixel area threshold,
\begin{equation}
T_t=\begin{cases}N_t^{-1}\sum_i\mathbf{1}[a_{i,t}<1024],&N_t>0,\\0,&N_t=0.\end{cases}
\end{equation}
The edge cue is the Canny edge fraction normalized by the historical edge norm of 0.14 and clipped to one. The night cue is the indicator that mean grayscale brightness is below 80. The blur cue is the indicator that Laplacian variance is below 180. These cues were selected because they are cheap relative to detector inference and correspond to different failure modes: density, target scale, image structure, illumination, and sharpness.

The historical raw SCI is
\begin{equation}
S_t=0.30C_t+0.30T_t+0.20E_t+0.10N_t+0.05B_t.
\end{equation}
The coefficients sum to 0.95, not one. This is an intentional property of the historical implementation and is retained in the description rather than silently renormalized. Analyses are performed every ten frames and the controller uses the mean of the most recent seven analysis values. Between analysis steps the operating point is held.

\subsection{Smart Calibrator}
The original confidence mapping is a decreasing function of scene complexity. Before clipping and the documented scene correction, it is
\begin{equation}
c_t=0.245-0.050S_t.
\end{equation}
For a scene classified as crowded, tiny, or night, the code applies a $-0.012$ correction and clips the result to $[0.19,0.28]$. The correction preserves more low-score observations when the scene cue indicates that a simple scalar index may underestimate difficulty.

The NMS mapping is similarly
\begin{equation}
\eta_t=0.490-0.050S_t.
\end{equation}
Blur applies a $-0.012$ correction, and the result is clipped to $[0.40,0.52]$. This adjustment changes how aggressively overlapping candidates are suppressed and is therefore detector-side, not a change to the tracker's association gate.

Resolution is selected by a discrete rule. The high-complexity action uses 832 pixels when $S_t>0.60$ or $T_t>0.50$. The medium action uses 736 pixels when $S_t>0.35$ or the scene label is crowded or tiny. The remaining state uses 640 pixels. The three controls are therefore adapted together, but they are not assumed to have identical monotonic effects on tracking quality.

The execution order is compactly: initialize the seven-entry SCI history and
previous tracker boxes; at the first frame and every tenth frame compute the
five cues from the current image and prior boxes, update the moving mean, and
apply the Smart Calibrator; then run the detector at the held operating point,
update the fixed tracker, and store its boxes for the next analysis step. The
full historical pseudocode is retained in Online Resource 1, Sec.~S3.
The architecture is causal but not necessarily optimal: a held state may lag a
scene transition, and a small-object action may increase association clutter.
Those are empirical questions addressed by the ablation and evolution evidence.

\subsection{Execution order and implementation semantics}
The temporal order is part of the method. At the beginning of a frame, the
controller has the previous stored tracker boxes and the current image; it
computes crowd and tiny-object statistics from the previous output and edge,
brightness, and blur from the current image. Only after detection and
association are the new boxes stored. Using current tracker output for its own
detector call would introduce a within-frame feedback loop and change the
historical causal design.

The image statistics use the reduced analysis image rather than a second
detector. Edge density is normalized, brightness is a mean grayscale statistic,
and blur is a Laplacian-variance statistic. An empty previous output gives a
zero tiny ratio, while crowd saturates at one after thirty boxes, making the
first decision well-defined with empty history.

Scene labels are computed alongside the scalar score: brightness determines
night, Laplacian variance determines blur, the area ratio determines tiny, and
density or edge conditions determine crowded. These labels permit the
deterministic confidence and NMS corrections that the scalar alone would not
express; no correction is learned on the final test split.

The seven-entry moving mean prevents one noisy analysis frame from changing the
detector immediately, at the cost of response delay. The stride and window are
historical design values, not universal optima. A reimplementation should
preserve the initialization, empty-box convention, stride, window, scene
thresholds, corrections, clipping ranges, and resolution order; implementation
details and the full trace are in Online Resource 1, Sec.~S3.
